\section{Related Work}

\subsection{Dense and Multi-Vector Information Retrieval}
Dense retrieval utilizing dual-encoder (bi-encoder) architectures has transformed modern information retrieval~\cite{reimers2019sentencebert,karpukhin2020dense}. Large-context multilingual embedding models such as BGE-M3~\cite{chen2024bge} support extended sequence windows (up to $8{,}192$ tokens) and multi-granularity signals. Nevertheless, for long, semi-structured documents like resumes, compressing heterogeneous sections into one single vector causes semantic dilution, while lightweight CPU-friendly encoders like \texttt{all-MiniLM-L6-v2} truncate documents beyond $256$ tokens. Token-level late interaction architectures such as ColBERT~\cite{khattab2020colbert} preserve fine-grained matching via MaxSim over token embeddings, but introduce prohibitive index storage expansion ($30\times$--$100\times$) and matrix compute overhead on CPUs. Recently, \v{C}eovi\'{c} et al.~\cite{ceovic2026multiview} proposed multi-view candidate--job matching using contrastive learning and LLM-based paraphrase augmentation. In contrast, our section-aware approach operates at the coarse functional section level ($M \le 4$ vectors per document), avoiding token-level index bloat and generative augmentation costs while preventing truncation and enabling per-section interpretability.

\subsection{Candidate--Job Matching and Skill Taxonomies}
Person--job fit modeling has progressed from bag-of-words retrieval to deep representation learning and graph neural architectures~\cite{zhu2018person,bian2020learning,zhang2021personalized}. Industrial hiring benchmarks highlight the complementary strengths of dense embeddings and lexical BM25~\cite{robertson2009probabilistic} fusion via methods such as Reciprocal Rank Fusion (RRF)~\cite{cormack2009reciprocal}, as recently evaluated in enterprise semantic candidate matching by Yashwant et al.~\cite{yashwant2026semantic}. However, conventional set-based skill matching relies on symmetric Jaccard similarity, which unfairly penalizes senior candidates who possess skills beyond a job's immediate checklist. Our work addresses this via asymmetric requirement coverage combined with a typed multi-hop skill ontology that distinguishes superset relationships, frameworks, and ecosystem companions while strictly guarding against false positive cross-ecosystem credit.

\subsection{Multi-Stage Reranking and Monotonic Learning-to-Rank}
High-precision retrieval pipelines frequently employ a cascaded two-stage design~\cite{asadi2013effectiveness,nogueira2019passage}, pairing fast first-stage candidate generation with cross-encoder rerankers. Recent systems such as ConFit~v3~\cite{jasonyux2026confitv3} apply reinforcement learning and listwise distillation over open-weights LLMs (Qwen3-8B/32B) to rerank lengthy resumes. While effective on GPU clusters, cross-encoders and LLM rerankers suffer severe CPU latency penalties and domain-shift degradation when trained on short passage QA benchmarks like MS~MARCO. Moreover, unconstrained neural rerankers lack structural monotonicity guarantees. We directly test this compute--accuracy trade-off against single-stage Calibrated Hybrid scoring and further evaluate monotonicity-constrained Learning-to-Rank (simplex-projected pairwise linear LTR and monotone-constrained LambdaMART) over multi-granularity section, lexical, and taxonomy signals.
